% Source: https://artofproblemsolving.com/community/c3h1337628_a_bit_of_bmmt
% Stable identity: aops-topic-1337628-q11
\textbf{Direct count.} Every subcommittee must contain at least one woman and at most two (three women would be all female, and with no woman it would be all male). With four women and three subcommittees, the women are split $2+1+1$. Choose the pair of women who serve together: $\dbinom42=6$ ways. The three subcommittees are now distinguished by their women. The pair gets one man ($5$ ways), and the remaining four men are split $2+2$ between the two single women ($\dbinom42=6$ ways). Each subcommittee then has one or two women, so none is single-gender. Total:
\[
6\cdot5\cdot6=\boxed{180}.
\]

\textbf{The student's argument is correct.} Fixing one person and choosing that person's two partners, then repeating for the next unplaced person, gives $\dbinom82\dbinom52=28\cdot10=280$ unlabelled splits; this equals $\dfrac{9!}{(3!)^3\,3!}=280$, so it is not ``too small''. An all-female subcommittee can be chosen in $4$ ways and the other six people split in $10$ ways ($40$); an all-male one in $\dbinom53=10$ ways and the rest in $10$ ways ($100$; two all-male subcommittees would need six men). As the student notes, an all-female subcommittee leaves five men and one woman, so one of the remaining subcommittees is all male: the overlap is all $40$. Hence $280-100=180$.

The key's $360$ is twice the correct count; it cannot be right for unlabelled subcommittees (and labelled subcommittees would give $6\cdot180=1080$).
